%=================================================================
% CSE 4345 — Big Data Analytics
% Week 1 Lecture Notes: Digital Data Landscape, Introduction to
%   Big Data, and the 3Vs Framework
% Department of CSE, Premier University
%=================================================================

\documentclass[12pt,a4paper]{article}

%---------- Core packages ----------
\usepackage[utf8]{inputenc}
\usepackage[T1]{fontenc}
\usepackage[margin=2.5cm, top=2.8cm, bottom=2.8cm]{geometry}
\usepackage{xcolor}
\usepackage{booktabs}
\usepackage{amsmath}
\usepackage{graphicx}
\usepackage{enumitem}
\usepackage{array}
\usepackage{fancyhdr}
\usepackage{titlesec}
\usepackage{parskip}
\usepackage{hyperref}
\usepackage{multicol}

%---------- Premier University Colors ----------
\definecolor{PUNavy}{RGB}{10,36,99}
\definecolor{PUGold}{RGB}{197,151,37}
\definecolor{PULightBlue}{RGB}{210,228,255}
\definecolor{PUTeal}{RGB}{0,100,100}
\definecolor{PUGray}{RGB}{90,90,90}
\definecolor{PURed}{RGB}{160,30,30}
\definecolor{PUGreen}{RGB}{30,120,60}

%---------- Hyperref ----------
\hypersetup{
  colorlinks=true,
  linkcolor=PUNavy,
  urlcolor=PUTeal,
  citecolor=PUGray
}

%---------- Section styling ----------
\titleformat{\section}
  {\large\bfseries\color{PUNavy}}
  {\thesection.}{0.6em}{}
  [\vspace{-2pt}\color{PUGold}\rule{\linewidth}{1.2pt}]

\titleformat{\subsection}
  {\normalsize\bfseries\color{PUNavy!85}}
  {\thesubsection}{0.6em}{}

\titleformat{\subsubsection}
  {\normalsize\bfseries\color{PUTeal}}
  {\thesubsubsection}{0.5em}{}

%---------- Page header / footer ----------
\pagestyle{fancy}
\fancyhf{}
\renewcommand{\headrulewidth}{0.6pt}
\renewcommand{\headrule}{\hbox to\headwidth{\color{PUNavy}\leaders\hrule height \headrulewidth\hfill}}
\lhead{\textcolor{PUNavy}{\small\textbf{CSE 4345 — Big Data Analytics}}}
\rhead{\textcolor{PUGray}{\small Week 1 Lecture Notes}}
\cfoot{\textcolor{PUGray}{\small Department of CSE, Premier University \textbar\ Page \thepage}}

%---------- Reusable box commands ----------
% Definition box
\newcommand{\defbox}[2]{%
  \noindent
  \colorbox{PUNavy!12}{%
    \begin{minipage}{\dimexpr\linewidth-2\fboxsep}
      \vspace{4pt}
      \textbf{\textcolor{PUNavy}{#1}}\par\vspace{2pt}
      #2
      \vspace{4pt}
    \end{minipage}%
  }\par\vspace{6pt}%
}

% Key insight box
\newcommand{\insightbox}[1]{%
  \noindent
  \colorbox{PUGold!18}{%
    \begin{minipage}{\dimexpr\linewidth-2\fboxsep}
      \vspace{4pt}
      \textcolor{PUGold!70!black}{\textbf{Key Insight}} \quad #1
      \vspace{4pt}
    \end{minipage}%
  }\par\vspace{6pt}%
}

% Example box
\newcommand{\examplebox}[2]{%
  \noindent
  \colorbox{PUTeal!10}{%
    \begin{minipage}{\dimexpr\linewidth-2\fboxsep}
      \vspace{4pt}
      \textbf{\textcolor{PUTeal}{Example: #1}}\par\vspace{2pt}
      #2
      \vspace{4pt}
    \end{minipage}%
  }\par\vspace{6pt}%
}

% Caution box
\newcommand{\cautionbox}[1]{%
  \noindent
  \colorbox{PURed!10}{%
    \begin{minipage}{\dimexpr\linewidth-2\fboxsep}
      \vspace{4pt}
      \textcolor{PURed}{\textbf{Common Misconception:}} \; #1
      \vspace{4pt}
    \end{minipage}%
  }\par\vspace{6pt}%
}

%=================================================================
\begin{document}
%=================================================================

%---------- Cover block ----------
\noindent
\colorbox{PUNavy}{%
  \begin{minipage}{\linewidth}
    \vspace{10pt}
    \begin{center}
      {\Large\bfseries\textcolor{white}{CSE 4345: Big Data Analytics}}\\[4pt]
      {\large\textcolor{PUGold}{Week 1 Lecture Notes}}\\[3pt]
      {\normalsize\textcolor{white}{Digital Data Landscape, Introduction to Big Data,
        and the 3Vs Framework}}\\[6pt]
      {\small\textcolor{PULightBlue}{Department of Computer Science \& Engineering
        \textbar\ Premier University, Chittagong}}\\[4pt]
      {\small\textcolor{PULightBlue}{\textbf{CLO:} CLO1 — Understand big data
        characteristics \quad\textbf{PLO:} PLO(a)}}
    \end{center}
    \vspace{8pt}
  \end{minipage}%
}

\vspace{14pt}

%---------- Learning Objectives ----------
\noindent
\colorbox{PULightBlue!70}{%
  \begin{minipage}{\dimexpr\linewidth-2\fboxsep}
    \vspace{5pt}
    \textbf{\textcolor{PUNavy}{After studying these notes you will be able to:}}
    \begin{itemize}[nosep, leftmargin=1.5em, after=\vspace{4pt}]
      \item Define Big Data using the 3Vs (and 5Vs) framework with real-world examples
      \item Distinguish between structured, semi-structured, and unstructured data
      \item Explain why traditional relational databases cannot handle Big Data workloads
      \item Differentiate between vertical and horizontal scaling
      \item State and apply the CAP theorem to distributed database design
      \item Compare batch processing and stream processing and identify suitable use cases for each
    \end{itemize}
  \end{minipage}%
}

\vspace{8pt}

%=================================================================
\section{The Data Revolution: Why Now?}
%=================================================================

Every minute of 2024, the world generates staggering volumes of data: 500 hours of video are uploaded to YouTube, 695{,}000 Instagram Stories are shared, 16 million text messages are sent, and over 6{,}000 tweets are posted. This scale is not just large — it is qualitatively different from anything that existed a generation ago.

\subsection{Historical Perspective}

Data has always existed, but the \textbf{rate of generation and the cost of storage} have changed beyond recognition:

\begin{center}
\renewcommand{\arraystretch}{1.4}
\begin{tabular}{p{2.0cm}p{4.2cm}p{4.2cm}p{2.8cm}}
  \toprule
  \textbf{Era} & \textbf{Data Sources} & \textbf{Scale} & \textbf{Dominant Tool} \\
  \midrule
  1960s--80s & Census, financial records & Megabytes & Mainframes \\
  1990s--2000s & Web clicks, e-commerce logs & Gigabytes -- Terabytes & RDBMS (Oracle, MySQL) \\
  2010s & Social media, smartphones, IoT & Petabytes -- Exabytes & Hadoop, NoSQL \\
  2020s & AI training data, genomics, autonomous vehicles & Zettabytes & Data lakes, lakehouse, streaming \\
  \bottomrule
\end{tabular}
\end{center}

\insightbox{By 2025, the IDC \textit{Global DataSphere} report forecasts the world's total stored data will reach \textbf{175 zettabytes} (175 $\times$ 10$^{21}$ bytes). If you stored this on DVDs, the stack would circle Earth 222 times.}

\subsection{Why Organizations Care}

The McKinsey Global Institute (2011) report \textit{``Big Data: The Next Frontier for Innovation, Competition, and Productivity''} (Manyika et al.) quantified the business case:
\begin{itemize}[leftmargin=1.8em]
  \item Retailers applying big data analytics could increase operating margins by \textbf{more than 60\%}
  \item The US healthcare system could unlock \textbf{over \$300 billion} in annual value through clinical analytics
  \item Big data is now a \textbf{factor of production} — alongside capital, land, and labor — that creates economic surplus
\end{itemize}

For Bangladesh specifically: the fintech sector (bKash, Nagad) processes hundreds of millions of transactions daily; DGHS manages health records for 170 million people; the RMG sector tracks logistics for thousands of factories. Each of these is a big data problem demanding techniques beyond what a standard relational database can provide.

%=================================================================
\section{Defining Big Data: The 3Vs Framework}
%=================================================================

The term ``Big Data'' is often used loosely. Its most widely accepted technical definition comes from a 2001 research note by \textbf{Doug Laney} at the META Group (later Gartner):

\defbox{Definition (Laney, 2001)}{Big Data is characterized by three dimensions:\newline
\textbf{Volume} (scale of data), \textbf{Velocity} (speed of generation and processing), and \textbf{Variety} (diversity of data types). Data is ``big'' when its Volume, Velocity, or Variety exceeds the capacity of conventional database tools to capture, manage, and process it within a tolerable elapsed time.}

These three dimensions are not independent — a system often faces all three simultaneously.

\subsection{Volume: Scale of Data}

Volume refers to the \textbf{sheer quantity} of data. The challenge is not just storage cost, but the inability of single-machine algorithms (which assume all data fits in RAM or even on one disk) to process it.

\begin{center}
\renewcommand{\arraystretch}{1.35}
\begin{tabular}{p{3.0cm}p{5.5cm}p{4.0cm}}
  \toprule
  \textbf{Organization} & \textbf{Data Volume} & \textbf{System Used} \\
  \midrule
  Facebook (Meta) & 500 TB ingested per day; 300 PB total & Apache Hive on Hadoop \\
  Google & 20 PB processed per day by MapReduce & GFS + Bigtable \\
  CERN (LHC) & 15 PB generated per year from particle collisions & Custom distributed stores \\
  bKash (Bangladesh) & $>$500 million transactions per year & Cloud + Hadoop stack \\
  NID Database (Bangladesh) & $\sim$110 million citizen records & Government data center \\
  \bottomrule
\end{tabular}
\end{center}

\textbf{Unit progression:} 1 KB $\to$ 1 MB $\to$ 1 GB $\to$ 1 TB $\to$ 1 PB $\to$ 1 EB $\to$ 1 ZB. Each step is a factor of 1,024. A single genome sequence is $\sim$200 GB; sequencing all humans on Earth would produce $\sim$1.5 ZB.

\subsection{Velocity: Speed of Generation and Processing}

Velocity refers to \textbf{how fast data arrives} and, crucially, \textbf{how quickly it must be processed} to be useful.

\begin{center}
\renewcommand{\arraystretch}{1.35}
\begin{tabular}{p{4.5cm}p{3.5cm}p{4.5cm}}
  \toprule
  \textbf{Source} & \textbf{Generation Rate} & \textbf{Processing Requirement} \\
  \midrule
  Twitter / X & 6{,}000 tweets/second & Near real-time trend detection \\
  Credit card transactions & 45{,}000 tx/second (Visa peak) & Fraud detection in $<$100 ms \\
  IoT sensors (factory floor) & Millions of readings/second & Predictive maintenance \\
  Pathao GPS pings & Every 3 seconds per active driver & Real-time routing \\
  Stock exchange ticks & $>$1 million/second & Algorithmic trading in microseconds \\
  \bottomrule
\end{tabular}
\end{center}

\insightbox{Velocity is why batch processing (e.g., running a nightly Hadoop job) is insufficient for some use cases. A fraud detection system that reports suspicious transactions the \textit{next morning} is useless — the money is already gone.}

\subsection{Variety: Diversity of Data Types}

Variety refers to the \textbf{diversity of formats and structures} in which data arrives.

\begin{center}
\renewcommand{\arraystretch}{1.4}
\begin{tabular}{p{2.8cm}p{5.0cm}p{4.8cm}}
  \toprule
  \textbf{Type} & \textbf{Characteristics} & \textbf{Examples} \\
  \midrule
  \textbf{Structured} & Fixed schema, tabular, easy to query &
    SQL databases, CSV files, spreadsheets \\
  \textbf{Semi-structured} & Self-describing schema, flexible &
    JSON, XML, log files, email headers \\
  \textbf{Unstructured} & No predefined format, hard to parse &
    Images, video, audio, free-text documents, MRI scans \\
  \bottomrule
\end{tabular}
\end{center}

A key fact worth memorizing: \textbf{structured data represents only $\sim$20\% of all digital data}. The remaining $\sim$80\% is unstructured — images, videos, documents, social media posts. Traditional RDBMS tools were built for structured data only, which is why big data ecosystems (Hadoop, Spark) extend to semi-structured and unstructured formats.

\subsection{The 5Vs Extension}

Later scholars and Gartner analysts extended Laney's 3Vs to acknowledge two additional dimensions:

\begin{center}
\renewcommand{\arraystretch}{1.4}
\begin{tabular}{p{2.4cm}p{6.0cm}p{4.2cm}}
  \toprule
  \textbf{V} & \textbf{Meaning} & \textbf{Example Challenge} \\
  \midrule
  \textbf{Volume}   & Quantity of data & 500 TB/day at Facebook \\
  \textbf{Velocity} & Speed of generation and required processing & Fraud detection in $<$100 ms \\
  \textbf{Variety}  & Diversity of formats and structures & Mixing JSON logs with MRI images \\
  \textbf{Veracity} & Uncertainty and quality of data & Sensor dropout, OCR errors, bots \\
  \textbf{Value}    & Business utility extracted from data & Revenue from predictive analytics \\
  \bottomrule
\end{tabular}
\end{center}

\defbox{Gartner Definition (Beyer \& Laney, 2012)}{``Big data is high-volume, high-velocity and/or high-variety information assets that demand cost-effective, innovative forms of information processing that enable enhanced insight, decision making, and process automation.''}

\cautionbox{``Big Data'' does not simply mean ``large dataset.'' A 10 TB single flat file with no velocity or variety requirements can be handled by a well-tuned PostgreSQL instance. Bigness is relative to your processing capacity, not an absolute number.}

%=================================================================
\section{Types of Digital Data}
%=================================================================

Understanding data types is foundational because each type requires \textbf{different storage systems, query engines, and analytics approaches}.

\subsection{Structured Data}

Structured data follows a \textbf{predefined schema}: every record has the same fields, with the same data types, in the same order. It is the native format of relational database management systems (RDBMS).

\textbf{Characteristics:}
\begin{itemize}[leftmargin=1.8em]
  \item Schema is enforced at write time (schema-on-write)
  \item Easily queried with SQL: \texttt{SELECT}, \texttt{JOIN}, \texttt{GROUP BY}
  \item Optimized for OLTP (Online Transaction Processing): fast inserts, point queries
  \item Represents approximately \textbf{20\%} of all digital data
\end{itemize}

\textbf{Examples:} bank transaction tables (account number, amount, timestamp, branch code), hospital patient records (ID, name, DOB, diagnosis code), bKash wallet balance tables.

\subsection{Semi-Structured Data}

Semi-structured data has \textbf{some organizational structure} but does not conform to a rigid tabular schema. It is self-describing through tags, keys, or markers.

\textbf{Characteristics:}
\begin{itemize}[leftmargin=1.8em]
  \item Schema is flexible and can vary between records (nested, repeated fields)
  \item No strict schema enforcement at ingestion
  \item Queryable via path expressions (JSON path, XPath) or schema-on-read tools (Hive, Spark SQL)
\end{itemize}

\textbf{Examples:}
\begin{itemize}[leftmargin=1.8em]
  \item \textbf{JSON}: \texttt{\{"user": "Alice", "amount": 1500, "tags": ["food","online"]\}} — common in web APIs
  \item \textbf{XML}: used in HL7 FHIR health record exchange between hospitals
  \item \textbf{Apache web logs}: \texttt{192.168.1.1 - - [12/Jun/2024:10:34:01] "GET /api/balance HTTP/1.1" 200 345}
  \item \textbf{Email headers}: structured fields (From, To, Subject) + unstructured body
\end{itemize}

\subsection{Unstructured Data}

Unstructured data has \textbf{no predefined format or schema}. It is the hardest to store, query, and analyse — and yet it is the most abundant form of data ($\sim$80\% of all digital data).

\textbf{Characteristics:}
\begin{itemize}[leftmargin=1.8em]
  \item Cannot be stored meaningfully in relational tables without heavy preprocessing
  \item Requires machine learning (computer vision, NLP, speech recognition) to extract meaning
  \item Often stored as files in distributed file systems (HDFS, S3) and object stores
\end{itemize}

\textbf{Examples:} CCTV video footage, medical imaging (X-ray, MRI, CT scan), satellite imagery (used by Bangladesh Meteorological Department), WhatsApp voice messages, court documents, social media posts in Bangla.

\examplebox{Healthcare Domain}{A hospital generates all three types simultaneously. The \textit{patient demographics table} (name, DOB, NID, blood type) is structured. The \textit{doctor's notes} in the Electronic Health Record are unstructured free text. The \textit{HL7 FHIR discharge summary} transmitted to another hospital is semi-structured XML. Big data platforms must ingest and correlate all three.}

%=================================================================
\section{Why Traditional Databases Cannot Handle Big Data}
%=================================================================

Traditional RDBMS (MySQL, Oracle, PostgreSQL) were designed for a world where:
\begin{enumerate}[leftmargin=2em]
  \item Data fits on a single server (or a small cluster with shared storage)
  \item Queries are transactional and require ACID guarantees
  \item The schema is known in advance and changes rarely
\end{enumerate}

When data exceeds these assumptions, three critical breakdowns occur:

\begin{center}
\renewcommand{\arraystretch}{1.45}
\begin{tabular}{p{3.2cm}p{4.5cm}p{5.0cm}}
  \toprule
  \textbf{Limitation} & \textbf{RDBMS Behavior} & \textbf{Big Data Consequence} \\
  \midrule
  Single-node capacity & Maximum data = one server's disk (a few TB) & Cannot store 500 TB; requires distributed storage \\
  Query latency & Full-table scan on 1 TB = hours & Analytics teams wait unacceptably long \\
  Schema rigidity & ALTER TABLE on 1 TB table = downtime & Cannot evolve schema as requirements change \\
  Variety blindness & No support for images, JSON trees, streams & Cannot analyse 80\% of enterprise data \\
  Concurrency under load & Write locks slow throughput & Cannot handle 45{,}000 tx/second ingestion \\
  \bottomrule
\end{tabular}
\end{center}

\subsection{Vertical vs.\ Horizontal Scaling}

The classical database solution to capacity problems is \textbf{vertical scaling} (scale up): buy a bigger, faster server with more RAM, faster CPU, and larger disks.

\begin{center}
\renewcommand{\arraystretch}{1.4}
\begin{tabular}{p{3.0cm}p{5.0cm}p{4.8cm}}
  \toprule
  & \textbf{Vertical Scaling (Scale Up)} & \textbf{Horizontal Scaling (Scale Out)} \\
  \midrule
  Approach & Add more RAM/CPU/disk to one machine & Add more machines to a cluster \\
  Cost & Exponential — high-end servers cost 10--100$\times$ more for 2$\times$ capacity & Linear — commodity servers (1{,}000 USD each) \\
  Fault tolerance & Single point of failure & Redundancy built across nodes \\
  Upper limit & Physical hardware limit ($\sim$few TB RAM, few PB disk) & Effectively unlimited (Google: millions of nodes) \\
  RDBMS fit & Natural fit; RDBMS was designed for this & Requires redesign of data and query model \\
  Big data approach & Not used & \textbf{Core design principle of Hadoop, Spark, Cassandra} \\
  \bottomrule
\end{tabular}
\end{center}

\insightbox{Hadoop's entire design philosophy is ``bring computation to data, not data to computation.'' Rather than moving petabytes to a central processing unit, Hadoop sends the processing code to each node where the data already lives. This \textbf{data locality} principle is the reason horizontal scaling works for big data analytics.}

%=================================================================
\section{The CAP Theorem}
%=================================================================

As soon as data is distributed across multiple machines (horizontal scaling), a fundamental question arises: \textit{what guarantees can the system make when the network between nodes fails?} Eric Brewer formally answered this in 2000.

\defbox{CAP Theorem (Brewer, 2000; formalized by Gilbert \& Lynch, 2002)}{A distributed data store can simultaneously guarantee \textbf{at most two} of the following three properties:\vspace{4pt}
\begin{itemize}[nosep, leftmargin=1.5em, after=\vspace{2pt}]
  \item \textbf{Consistency (C):} Every read receives the most recent write (or an error) — all nodes see the same data at the same time
  \item \textbf{Availability (A):} Every request receives a (non-error) response, though it may not contain the most recent data
  \item \textbf{Partition Tolerance (P):} The system continues to operate even when network messages between nodes are lost or delayed (a network partition)
\end{itemize}}

\subsection{Understanding the Three Properties}

\textbf{Consistency} means that if you write a value to node A, and immediately read it from node B (a different server), you get the written value — not an old cached copy. This is what you expect from a traditional bank transfer: if you deposit \$100, your balance on any ATM worldwide should reflect that immediately.

\textbf{Availability} means every request gets a response. The system never rejects a query with ``I cannot answer right now.'' Even if some nodes are down, the remaining nodes serve requests.

\textbf{Partition Tolerance} means the system keeps working even when network communication between nodes breaks. In a real distributed system — spanning data centers, countries, or even just two racks in one building — network partitions are \textbf{not hypothetical}. They happen regularly due to cable cuts, switch failures, and routing errors.

\subsection{The Critical Implication}

Since network partitions \textbf{will occur} in any real distributed system, Partition Tolerance (P) is non-negotiable for large-scale systems. This reduces the real choice to:
\begin{itemize}[leftmargin=2em]
  \item \textbf{CP systems} (Consistency + Partition Tolerance): When a partition occurs, reject some requests to ensure the data that \textit{is} returned is always correct. \textbf{Examples:} HBase, ZooKeeper, MongoDB (default config), Google Spanner
  \item \textbf{AP systems} (Availability + Partition Tolerance): When a partition occurs, continue serving requests even if responses may return slightly stale data. \textbf{Examples:} Cassandra, DynamoDB, CouchDB
\end{itemize}

\begin{center}
\renewcommand{\arraystretch}{1.45}
\begin{tabular}{p{3.8cm}p{2.4cm}p{5.8cm}}
  \toprule
  \textbf{System} & \textbf{CAP Choice} & \textbf{Reason / Trade-off} \\
  \midrule
  Traditional RDBMS (MySQL) & CA & Runs on one node; partitions don't apply \\
  HBase & CP & Prefers to reject requests than serve stale data (banking) \\
  Apache Cassandra & AP & Eventual consistency; prefers uptime (social feed, IoT) \\
  Amazon DynamoDB & AP (tunable) & Defaults to eventual consistency; strong consistency optional \\
  Google Spanner & CP (global) & Uses GPS + atomic clocks to achieve strong consistency at scale \\
  ZooKeeper & CP & Coordination service: consistency is paramount \\
  \bottomrule
\end{tabular}
\end{center}

\examplebox{bKash Mobile Banking}{bKash processes financial transfers. It uses a \textbf{CP-leaning} design: if the system cannot guarantee that a transfer is atomically applied across all replicas, it \textit{refuses} the transaction rather than risking a double-spend or balance discrepancy. Availability is sacrificed briefly (you may get a ``try again'' error) to protect Consistency. Contrast this with a social media feed (like Facebook's News Feed), which uses an \textbf{AP design}: it is acceptable to see a post 2 seconds late; it is not acceptable for the feed to go down.}

\cautionbox{The CAP theorem does not mean ``choose two, ignore one completely.'' Modern systems (like Cassandra with tunable consistency, or Google Spanner) are designed to \textit{minimize} the cases where the trade-off bites, using techniques like quorum reads/writes and synchronized clocks. Think of CAP as describing behavior \textit{during a partition}, not normal operation.}

%=================================================================
\section{Batch Processing vs.\ Stream Processing}
%=================================================================

Once data is stored in a distributed system, there are two fundamentally different paradigms for processing it:

\defbox{Batch Processing}{Processing of data \textbf{collected over a period of time} as a single bounded job. The system reads all input, applies a computation, and produces output. The input dataset is complete and static during processing.\newline\textbf{Canonical system:} Apache Hadoop MapReduce.}

\defbox{Stream Processing}{Processing of data \textbf{as it arrives}, record by record or in micro-batches, over an unbounded (potentially infinite) data stream. The system produces output continuously, often within milliseconds of data arrival.\newline\textbf{Canonical systems:} Apache Kafka, Apache Flink, Spark Structured Streaming.}

\subsection{Side-by-Side Comparison}

\begin{center}
\renewcommand{\arraystretch}{1.5}
\begin{tabular}{p{3.5cm}p{4.5cm}p{4.5cm}}
  \toprule
  \textbf{Property} & \textbf{Batch Processing} & \textbf{Stream Processing} \\
  \midrule
  Data bounds & Bounded (finite dataset) & Unbounded (infinite stream) \\
  Latency & Minutes to hours & Milliseconds to seconds \\
  Throughput & Very high (optimized for bulk) & Lower per unit time \\
  Complexity & Lower (no state management) & Higher (windowing, watermarks, state) \\
  Fault tolerance & Re-run from input file & Checkpointing, replay from Kafka \\
  Output & Written when job completes & Continuous (updated results) \\
  Typical use case & Nightly ETL, monthly reports, model training & Fraud detection, live dashboards, alerting \\
  \textbf{Examples} & Hadoop MR, Hive, Spark batch & Kafka Streams, Flink, Spark Streaming \\
  \bottomrule
\end{tabular}
\end{center}

\subsection{Choosing the Right Paradigm}

The decision between batch and stream processing is driven by \textbf{how quickly the result must be available after the data arrives}:

\begin{itemize}[leftmargin=1.8em]
  \item \textbf{Use batch when:} results are needed hours or days after data collection (monthly sales reports, end-of-day bank reconciliation, weekly model retraining, DGHS annual disease statistics)
  \item \textbf{Use stream when:} results are needed within seconds of data arrival (bKash fraud alert, Pathao surge pricing when 100 drivers go offline simultaneously, hospital patient monitoring alarms, stock trading signals)
  \item \textbf{Use both (Lambda / Kappa architecture):} many production systems combine batch for historical accuracy and stream for low-latency approximations (covered in Week 8)
\end{itemize}

\examplebox{Ride-Sharing: Pathao in Bangladesh}{Pathao collects GPS pings from drivers every 3 seconds. \textbf{Stream processing} is used to update the live map of available drivers in real time (latency requirement: $<$1 second). \textbf{Batch processing} is used to compute weekly driver earnings reports, trip analytics, and demand heatmaps by neighbourhood (latency requirement: hours to days, but covering millions of historical records). Both pipelines run in parallel on the same underlying data.}

%=================================================================
\section{Summary and Key Takeaways}
%=================================================================

\begin{center}
\renewcommand{\arraystretch}{1.5}
\begin{tabular}{p{4.0cm}p{9.0cm}}
  \toprule
  \textbf{Concept} & \textbf{What to Remember} \\
  \midrule
  3Vs (Laney, 2001) & Volume, Velocity, Variety — the canonical definition of Big Data; extended to 5Vs with Veracity and Value \\
  Structured data & $\sim$20\% of all data; fixed schema, queryable with SQL; RDBMS territory \\
  Unstructured data & $\sim$80\% of all data; images, video, free text; requires ML to process \\
  Vertical scaling & Add more to one machine; expensive; has a hard limit \\
  Horizontal scaling & Add more machines; linear cost; the core principle of Hadoop/Spark \\
  CAP theorem & Distributed systems can guarantee only 2 of C, A, P simultaneously; since P is required, real choice is CP vs.\ AP \\
  CP systems & Consistency + Partition Tolerance; reject requests during partition (e.g., HBase, ZooKeeper) \\
  AP systems & Availability + Partition Tolerance; return possibly stale data during partition (e.g., Cassandra, DynamoDB) \\
  Batch processing & Bounded input, high throughput, minutes-to-hours latency; nightly ETL, reports \\
  Stream processing & Unbounded input, low latency (ms--s); fraud detection, live dashboards \\
  \bottomrule
\end{tabular}
\end{center}

\insightbox{The three Vs of Big Data (Volume, Velocity, Variety) define the \textit{problem}. Horizontal scaling, the CAP theorem, and batch/stream architectures define the \textit{engineering responses} to that problem. Everything in this course (Hadoop, Spark, Hive, Kafka) is an implementation of those responses.}

%=================================================================
\section{Self-Assessment}
%=================================================================

\subsection*{Multiple Choice Questions}

\noindent\textbf{Instructions:} Select the single best answer.

\begin{enumerate}[label=\textbf{Q\arabic*.}, leftmargin=2.5em, itemsep=10pt]

  \item Who first formally introduced the 3Vs framework for Big Data?
  \begin{enumerate}[label=(\alph*), leftmargin=2em, nosep, after=\vspace{4pt}]
    \item Jim Gray
    \item \textbf{Doug Laney} \textcolor{PUGreen}{\small $\leftarrow$ correct}
    \item Jeffrey Dean
    \item Erik Brynjolfsson
  \end{enumerate}

  \item Which of the following is an example of \textbf{unstructured data}?
  \begin{enumerate}[label=(\alph*), leftmargin=2em, nosep, after=\vspace{4pt}]
    \item A SQL database table of student records
    \item A CSV file of monthly sales figures
    \item \textbf{An MRI scan image} \textcolor{PUGreen}{\small $\leftarrow$ correct}
    \item A JSON array of product prices
  \end{enumerate}

  \item The CAP theorem states that a distributed system cannot simultaneously guarantee more than two of which three properties?
  \begin{enumerate}[label=(\alph*), leftmargin=2em, nosep, after=\vspace{4pt}]
    \item Cost, Availability, Performance
    \item \textbf{Consistency, Availability, Partition Tolerance} \textcolor{PUGreen}{\small $\leftarrow$ correct}
    \item Correctness, Atomicity, Persistence
    \item Concurrency, Accuracy, Partitioning
  \end{enumerate}

  \item Which ``V'' of Big Data refers to the \textbf{speed} at which data is generated and must be processed?
  \begin{enumerate}[label=(\alph*), leftmargin=2em, nosep, after=\vspace{4pt}]
    \item Volume
    \item Variety
    \item \textbf{Velocity} \textcolor{PUGreen}{\small $\leftarrow$ correct}
    \item Veracity
  \end{enumerate}

\end{enumerate}

\subsection*{True / False}

\noindent\textbf{Instructions:} State True or False and write a one-sentence justification.

\begin{center}
\renewcommand{\arraystretch}{1.55}
\begin{tabular}{p{9.5cm}c}
  \toprule
  \textbf{Statement} & \textbf{Answer} \\
  \midrule
  1.\ Structured data accounts for approximately 80\% of all digital data. &
  \textbf{\textcolor{PURed}{False}} \\
  2.\ The 3Vs framework was later extended to 5Vs by adding Veracity and Value. &
  \textbf{True} \\
  3.\ Horizontal scaling means adding more RAM or CPU to a single existing server. &
  \textbf{\textcolor{PURed}{False}} \\
  4.\ Batch processing is suitable for real-time fraud detection in payment systems. &
  \textbf{\textcolor{PURed}{False}} \\
  \bottomrule
\end{tabular}
\end{center}

\noindent\textbf{Justifications:}
\begin{enumerate}[leftmargin=2em, nosep, after=\vspace{6pt}]
  \item \textit{False} — Structured data is only $\sim$20\%; approximately 80\% of all digital data is unstructured (images, video, free text).
  \item \textit{True} — Gartner's Beyer \& Laney (2012) formally extended the 3Vs to 5Vs, adding Veracity (data quality uncertainty) and Value (business utility).
  \item \textit{False} — Horizontal scaling adds more machines to a cluster (scale out). Adding more hardware to a single server is vertical scaling (scale up).
  \item \textit{False} — Batch processing has latency of minutes to hours. Fraud detection requires sub-second decision-making, which requires stream processing.
\end{enumerate}

\subsection*{Descriptive Questions}
\noindent\textbf{Instructions:} Answer each question in 100--150 words.

\begin{enumerate}[label=\textbf{D\arabic*.}, leftmargin=2.5em, itemsep=8pt]
  \item Explain the 3Vs (and optionally 5Vs) of Big Data with one real-world example for each dimension.
  \item Differentiate between structured, semi-structured, and unstructured data with examples drawn from the healthcare domain (hospital, clinic, or public health).
  \item What is the CAP theorem? State each of the three properties in one sentence each. Why is Partition Tolerance considered non-negotiable for large-scale distributed systems?
\end{enumerate}

\subsection*{Analytical Questions}
\noindent\textbf{Instructions:} Answer in 200--300 words with step-by-step reasoning.

\begin{enumerate}[label=\textbf{A\arabic*.}, leftmargin=2.5em, itemsep=8pt]
  \item \textbf{5Vs Analysis of a Ride-Sharing Platform:}\\
  A ride-sharing company collects GPS pings every 3 seconds from 5 million active drivers. Each ping includes driver ID, latitude, longitude, speed, and battery level. Analyse this scenario using the full \textbf{5Vs} framework. For each V: (a) describe the specific challenge, and (b) propose a technical approach to address it. Which V poses the \textit{greatest} engineering challenge for this system and why?

  \item \textbf{Batch vs.\ Stream Processing Decision:}\\
  You are the lead data engineer for a large Bangladeshi bank. The bank needs to: (i) detect fraudulent transactions within 200 milliseconds, (ii) generate end-of-month customer account statements, (iii) retrain the fraud detection ML model weekly on the past 6 months of transactions, and (iv) power a live dashboard showing the total number of active sessions across all branches. For each of the four requirements, specify whether you would use batch processing, stream processing, or both, and justify your choice with reference to the latency, volume, and fault-tolerance properties discussed in this lecture.
\end{enumerate}

%=================================================================
\section*{Textbook References for Week 1}
%=================================================================

\begin{itemize}[leftmargin=2em, itemsep=6pt]
  \item \textbf{[SA]} Seema Acharya \& Subhashini Chellappan. \textit{Big Data and Analytics}. Wiley, 2015.\hfill\textbf{Chapter 1}\\
        \small\textit{Chapter 1 provides a foundational introduction to big data: the 3Vs, types of data, business drivers, and an overview of the Hadoop ecosystem. Primary reading for this week.}

  \item \textbf{[TE]} Thomas Erl, Wajid Khattak \& Paul Buhler. \textit{Big Data Fundamentals}. Prentice Hall, 2016.\hfill\textbf{Chapter 1}\\
        \small\textit{Chapter 1 defines big data from an enterprise architecture perspective, extending the 3Vs to include operational and enterprise big data types.}

  \item \textbf{[MMS]} Viktor Mayer-Schönberger \& Kenneth Cukier. \textit{Big Data: A Revolution}. Houghton Mifflin, 2013.\hfill\textbf{Chapter 1}\\
        \small\textit{Chapter 1 (``Now'') provides a narrative account of how big data is transforming society — excellent context for the ``why it matters'' framing of this week.}

  \item \textbf{[LRU]} Leskovec, Rajaraman \& Ullman. \textit{Mining of Massive Datasets}, 3rd ed. Cambridge Univ.\ Press, 2020.\hfill\textbf{Chapter 1}\\
        \small\textit{Chapter 1 frames the problem from a computational perspective: algorithms that must run on data too large for a single machine. Available free at \texttt{mmds.org}.}

  \item \textbf{Seminal Paper:} Doug Laney (2001). \textit{``3D Data Management: Controlling Data Volume, Velocity and Variety.''} META Group Research Note, 6 Feb 2001. [The original 3Vs paper — 2 pages, recommended reading]

  \item \textbf{Seminal Paper:} Eric Brewer (2000). \textit{``Towards Robust Distributed Systems.''} Keynote, ACM PODC 2000. [Original CAP conjecture; formalized by Gilbert \& Lynch, ACM SIGACT News, 2002]
\end{itemize}

\vspace{12pt}
\noindent\rule{\linewidth}{0.8pt}\\[4pt]
{\small\textcolor{PUGray}{CSE 4345 — Big Data Analytics \textbar\ Week 1 Lecture Notes \textbar\
Department of CSE, Premier University \textbar\ Instructor: rufaruqui@cu.ac.bd}}

\end{document}
%=================================================================
